\section{激活检查点（Activation Checkpointing）}

即使使用 FSDP 解决了权重存储问题，激活值的存储仍会成为瓶颈。以 Llama 3 405B 为例（126 层、维度 16384、序列长度 4096），前向传播的隐藏状态需要大量 GPU 内存。

\begin{figure}[H]
\centering
\includegraphics[width=0.95\textwidth]{slides-images/slide-024.jpg}
\caption{激活检查点原理：不存储所有中间激活，而是在反向传播时重新计算\protect\footnotemark}
\end{figure}
\footnotetext{Slides 第24页。}

\subsection{三种策略的计算-内存权衡}

对于 $n$ 层网络：

\begin{table}[H]
\centering
\begin{tabular}{lcc}
\toprule
\textbf{策略} & \textbf{计算} & \textbf{内存} \\
\midrule
存储所有激活 & $O(n)$ & $O(n)$ \\
不存储任何激活 & $O(n^2)$ & $O(1)$ \\
每 $c$ 层检查点一次 & $O(n^2/c)$ & $O(c)$ \\
$c = \sqrt{n}$（最优） & $O(n\sqrt{n})$ & $O(\sqrt{n})$ \\
\bottomrule
\end{tabular}
\caption{激活检查点的计算-内存权衡}
\end{table}

\begin{warningbox}{激活检查点会降低训练速度}
激活检查点通过在反向传播时重新计算中间激活来节省内存，但这会引入额外的计算开销。典型设置（$c = \sqrt{n}$）会增加约 $\sqrt{n}$ 倍的前向计算量。这``不太爽''，但它让我们能训练更大的模型。
\end{warningbox}

\subsection{本章小结}

激活检查点是一种用计算换内存的技术，在 FSDP 之上进一步扩展了可训练模型的规模。最优策略是每 $\sqrt{n}$ 层保存一次检查点，实现 $O(n\sqrt{n})$ 计算和 $O(\sqrt{n})$ 内存的平衡。

%% ========== 扩展策略 ==========

